\documentclass[a4paper,12pt]{report}

% Packages needed
\usepackage[utf8]{inputenc}
\usepackage[a4paper, margin=2.5cm]{geometry}
\usepackage{graphicx}
\usepackage{float}
\usepackage{booktabs}
\usepackage{listings}
\usepackage{xcolor}
\usepackage{hyperref}
\usepackage{capt-of}
\usepackage{titlesec}
\usepackage{wrapfig}
\usepackage{needspace}

% Reduce the large empty space before a chapter title (no "Chapter N" prefix shown)
\titleformat{\chapter}[display]
  {\normalfont\huge\bfseries}{}{0pt}{\Huge}
\titlespacing*{\chapter}{0pt}{-50pt}{20pt}

% Code listing style (Python)
\lstset{
  language=Python,
  basicstyle=\ttfamily\small,
  keywordstyle=\color{blue}\bfseries,
  commentstyle=\color{gray}\ttfamily,
  stringstyle=\color{red}\ttfamily,
  numbers=left,
  numberstyle=\tiny\color{gray},
  stepnumber=1,
  numbersep=5pt,
  frame=single,
  breaklines=true,
  breakatwhitespace=false,
  showstringspaces=false,
  tabsize=2,
  captionpos=b
}

\begin{document}

\setcounter{chapter}{2}
\setcounter{section}{0}
\chapter*{Image Processing Nodes}
\addcontentsline{toc}{chapter}{Image Processing Nodes}
\label{chap:image-processing-nodes}

The Image Processing Nodes provide a modular preprocessing pipeline for
computer vision applications. Each node performs a single image processing
task with configurable parameters and produces the processed image along
with execution metadata. This modular design allows different nodes to be
combined flexibly to build customized image processing workflows.

\section{Resize Node}
\label{sec:resize-node}

The Resize Node adjusts image dimensions to match the required input size
for subsequent processing stages. It supports multiple resizing strategies
while preserving image quality through different interpolation methods.

\noindent \textbf{Input:} Image, Target Size, Resize Mode, Interpolation Method \\
\textbf{Output:} Resized Image, Metadata \\
\textbf{Supported Modes:} Fit, Fill, Stretch

\begin{lstlisting}[caption={Core Algorithm}, label={lst:resize-node}]
if self.mode == "stretch":

    resized = img.resize(
        (target_w, target_h),
        self.interpolation
    )

elif self.keep_aspect_ratio:

    resized = self._resize_with_aspect(
        img,
        target_w,
        target_h
    )
\end{lstlisting}

\section{Crop Node}
\label{sec:crop-node}

The Crop Node extracts a Region of Interest (ROI) from the input image by
removing unnecessary surrounding areas. This operation reduces
computational complexity and enables subsequent nodes to focus on the
relevant image content.

\noindent \textbf{Input:} Image, Crop Coordinates / Crop Size \\
\textbf{Output:} Cropped Image, Metadata \\
\textbf{Supported Modes:} Rectangle Crop, Center Crop, ROI Crop, Percentage Crop

\begin{lstlisting}[caption={Core Algorithm}, label={lst:crop-node}]
cropped_image = image[
    y:y + height,
    x:x + width
]
\end{lstlisting}

\section{Transform Node}
\label{sec:transform-node}

The Transform Node performs geometric transformations on the input image
for preprocessing and data preparation. It also provides augmentation
capabilities by combining multiple transformation operations.

\noindent \textbf{Input:} Image, Transformation Parameters \\
\textbf{Output:} Transformed Image, Metadata \\
\textbf{Supported Modes:} Rotate, Rotate 90\textdegree, Flip, Affine Transformation, Augmentation

\begin{lstlisting}[caption={Core Algorithm}, label={lst:transform-node}]
if self.operation == "rotate":

    img, d = self._rotate(img)

elif self.operation == "flip":

    img, d = self._flip(img)

elif self.operation == "affine":

    img, d = self._affine(img)

elif self.operation == "augment":

    img, d = self._augment(img)
\end{lstlisting}

\section{Color Space Conversion Node}
\label{sec:color-space-conversion-node}

The Color Space Conversion Node converts images between different color
representations to satisfy the requirements of various computer vision
algorithms. Choosing an appropriate color space improves the performance
of many image analysis tasks.

\noindent \textbf{Input:} Image, Source Color Space, Destination Color Space \\
\textbf{Output:} Converted Image, Metadata \\
\textbf{Supported Modes:} RGB, BGR, Grayscale, HSV, LAB, YCrCb

\begin{lstlisting}[caption={Core Algorithm}, label={lst:color-space-conversion-node}]
result = cv2.cvtColor(
    image,
    conversion_code
)
\end{lstlisting}

\section{Normalize Node}
\label{sec:normalize-node}

The Normalize Node scales pixel values into standardized numerical ranges
suitable for machine learning and deep learning models. Different
normalization techniques are supported to satisfy the requirements of
different model architectures.

\noindent \textbf{Input:} Image, Normalization Mode \\
\textbf{Output:} Normalized Image, Metadata \\
\textbf{Supported Modes:} Min-Max, Mean-Std, Z-Score, Scale

\begin{lstlisting}[caption={Core Algorithm}, label={lst:normalize-node}]
if self.mode == "minmax":

    out, details = self._minmax(img)

elif self.mode == "meanstd":

    out, details = self._meanstd(img)

elif self.mode == "zscore":

    out, details = self._zscore(img)

elif self.mode == "scale":

    out, details = self._scale(img)
\end{lstlisting}

\section{Histogram Analysis Node}
\label{sec:histogram-analysis-node}

The Histogram Analysis Node analyzes the intensity distribution of
grayscale or color images by computing histogram values for different
channels. The generated statistics help evaluate image brightness,
contrast, and overall intensity distribution before further processing.

\noindent \textbf{Input:} Image, Analysis Mode \\
\textbf{Output:} Histogram Data, Metadata \\
\textbf{Supported Modes:} Gray, RGB, HSV, All

\begin{lstlisting}[caption={Core Algorithm}, label={lst:histogram-analysis-node}]
hist = cv2.calcHist(
    [image],
    [channel],
    None,
    [256],
    [0, 256]
)
\end{lstlisting}

\section{Image Enhancement Node}
\label{sec:image-enhancement-node}

The Image Enhancement Node improves the visual quality of images by
enhancing contrast, brightness, and illumination. Different enhancement
techniques are provided to improve image appearance while preserving
important visual details.

\noindent \textbf{Input:} Image, Enhancement Method, Enhancement Parameters \\
\textbf{Output:} Enhanced Image, Metadata \\
\textbf{Supported Modes:} CLAHE, Histogram Equalization, Gamma Correction, Manual Adjustment

\begin{lstlisting}[caption={Core Algorithm}, label={lst:image-enhancement-node}]
if safe_method == "clahe":

    clahe = cv2.createCLAHE(
        clipLimit=clip_limit,
        tileGridSize=tile_grid
    )

    lab = cv2.cvtColor(
        image,
        cv2.COLOR_BGR2LAB
    )

    l, a, b = cv2.split(lab)

    l = clahe.apply(l)
\end{lstlisting}

\section{Histogram Equalization Node}
\label{sec:histogram-equalization-node}

The Histogram Equalization Node enhances image contrast by redistributing
pixel intensity values. It supports both global and adaptive histogram
equalization techniques for handling images with different illumination
conditions.

\noindent \textbf{Input:} Image, Equalization Method \\
\textbf{Output:} Equalized Image, Metadata \\
\textbf{Supported Modes:} Global Histogram Equalization, CLAHE

\begin{lstlisting}[caption={Core Algorithm}, label={lst:histogram-equalization-node}]
if method == "clahe":

    clahe = cv2.createCLAHE(
        clipLimit=clip_limit,
        tileGridSize=tile_grid
    )

    result = clahe.apply(gray)

else:

    result = cv2.equalizeHist(gray)
\end{lstlisting}

\section{Blur \& Denoise Node}
\label{sec:blur-denoise-node}

The Blur \& Denoise Node reduces image noise while preserving important
structures and edges. Multiple filtering techniques are available to
achieve different levels of smoothing according to the application
requirements.

\noindent \textbf{Input:} Image, Filter Method, Kernel Size, Filter Parameters \\
\textbf{Output:} Denoised Image, Metadata \\
\textbf{Supported Modes:} Gaussian Blur, Median Blur, Box Blur, Bilateral Filter, Non-Local Means

\begin{lstlisting}[caption={Core Algorithm}, label={lst:blur-denoise-node}]
if safe_method == "gaussian":

    result = cv2.GaussianBlur(
        image,
        (k, k),
        sigmaX=sigma
    )

elif safe_method == "bilateral":

    result = cv2.bilateralFilter(
        image,
        d=k,
        sigmaColor=sigma,
        sigmaSpace=sigma
    )
\end{lstlisting}

\section{Thresholding Node}
\label{sec:thresholding-node}

The Thresholding Node converts grayscale images into binary images by
separating foreground objects from the background. It supports both
global and adaptive thresholding methods for different lighting
conditions.

\noindent \textbf{Input:} Image, Threshold Method, Threshold Parameters \\
\textbf{Output:} Binary Image, Metadata \\
\textbf{Supported Modes:} Global, Otsu, Adaptive Mean, Adaptive Gaussian

\begin{lstlisting}[caption={Core Algorithm}, label={lst:thresholding-node}]
if safe_method == "otsu":

    _, result = cv2.threshold(
        gray_image,
        0,
        max_val,
        thresh_type +
        cv2.THRESH_OTSU
    )

elif safe_method == "adaptive_gaussian":

    result = cv2.adaptiveThreshold(
        gray_image,
        max_val,
        cv2.ADAPTIVE_THRESH_GAUSSIAN_C,
        thresh_type,
        b_size,
        c_val
    )
\end{lstlisting}

\section{Morphology Node}
\label{sec:morphology-node}

The Morphology Node applies morphological operations to binary images in
order to remove noise, refine object boundaries, and improve segmentation
quality. The node supports different structuring element shapes and
multiple iterations.

\noindent \textbf{Input:} Binary Image, Morphological Operation, Kernel Shape, Kernel Size, Iterations \\
\textbf{Output:} Processed Image, Metadata \\
\textbf{Supported Modes:} Erosion, Dilation, Opening, Closing, Gradient

\begin{lstlisting}[caption={Core Algorithm}, label={lst:morphology-node}]
if safe_method == "erosion":

    result = cv2.erode(
        binary,
        kernel,
        iterations=iterations
    )

elif safe_method == "closing":

    result = cv2.morphologyEx(
        binary,
        cv2.MORPH_CLOSE,
        kernel,
        iterations=iterations
    )
\end{lstlisting}

\section{Edge Detection Node}
\label{sec:edge-detection-node}

The Edge Detection Node identifies object boundaries by detecting abrupt
intensity changes within the image. It provides gradient-based and
multi-stage edge detection algorithms that are widely used in segmentation
and feature extraction tasks.

\noindent \textbf{Input:} Image, Edge Detection Method, Detection Parameters \\
\textbf{Output:} Edge Image, Metadata \\
\textbf{Supported Modes:} Canny, Sobel

\begin{lstlisting}[caption={Core Algorithm}, label={lst:edge-detection-node}]
if safe_method == "canny":

    result = cv2.Canny(
        gray,
        threshold1,
        threshold2
    )

elif safe_method == "sobel":

    grad_x = cv2.Sobel(
        gray,
        cv2.CV_64F,
        1,
        0,
        ksize=ksize
    )

    grad_y = cv2.Sobel(
        gray,
        cv2.CV_64F,
        0,
        1,
        ksize=ksize
    )
\end{lstlisting}

\section{Connected Components Node}
\label{sec:connected-components-node}

The Connected Components Node labels connected foreground regions in a
binary image and extracts descriptive information for each detected
object. It can also visualize detected objects using bounding boxes and
labels.

\noindent \textbf{Input:} Binary Image, Connectivity, Area Constraints \\
\textbf{Output:} Labeled Image, Component Statistics, Metadata \\
\textbf{Supported Modes:} 4-Connectivity, 8-Connectivity

\begin{lstlisting}[caption={Core Algorithm}, label={lst:connected-components-node}]
num_labels, labels, stats, centroids = \
cv2.connectedComponentsWithStats(

    binary,

    connectivity=self.connectivity

)
\end{lstlisting}

\section{Contour \& Shape Detection Node}
\label{sec:contour-shape-detection-node}

The Contour \& Shape Detection Node detects the external contours of
objects within a binary image and calculates their geometric boundaries.
It is commonly used for object localization, measurement, and shape
analysis.

\noindent \textbf{Input:} Image, Minimum Area, Bounding Box Parameters \\
\textbf{Output:} Image with Detected Contours, Metadata \\
\textbf{Supported Modes:} External Contour Detection, Bounding Box Visualization

\begin{lstlisting}[caption={Core Algorithm}, label={lst:contour-shape-detection-node}]
contours, _ = cv2.findContours(

    thresh,

    cv2.RETR_EXTERNAL,

    cv2.CHAIN_APPROX_SIMPLE

)

for cnt in contours:

    area = cv2.contourArea(cnt)

    if area > min_area:

        x, y, w, h = cv2.boundingRect(cnt)
\end{lstlisting}

\section{Feature Extraction Node}
\label{sec:feature-extraction-node}

The Feature Extraction Node detects distinctive keypoints and computes
feature descriptors that uniquely represent local image structures. These
descriptors are later used in image matching, registration, and object
recognition.

\noindent \textbf{Input:} Image, Feature Extraction Method, Maximum Features \\
\textbf{Output:} Keypoints, Descriptors, Metadata \\
\textbf{Supported Modes:} ORB, BRISK, AKAZE

\begin{lstlisting}[caption={Core Algorithm}, label={lst:feature-extraction-node}]
if method == "orb":

    detector = cv2.ORB_create(

        nfeatures=max_features

    )

keypoints, descriptors = detector.detectAndCompute(

    gray,

    None

)
\end{lstlisting}

\section{Feature Matching Node}
\label{sec:feature-matching-node}

The Feature Matching Node compares feature descriptors extracted from two
different images to identify corresponding keypoints. It supports
brute-force matching and provides configurable filtering strategies to
retain only reliable matches.

\noindent \textbf{Input:} Two Images, Keypoints, Descriptors, Matching Parameters \\
\textbf{Output:} Matched Features, Matching Statistics, Metadata \\
\textbf{Supported Modes:} Brute Force (BF), Cross Check, Top-N Matching, Distance Threshold Filtering

\begin{lstlisting}[caption={Core Algorithm}, label={lst:feature-matching-node}]
matcher = cv2.BFMatcher(

    norm_type,

    crossCheck=cross_check

)

matches = matcher.match(

    descriptors1,

    descriptors2

)

matches = sorted(

    matches,

    key=lambda x: x.distance

)
\end{lstlisting}

\end{document}
